\section{引言}
假设两个模拟系统预测某项活动的到场人数相同。第一个系统中，几乎所有人都收到通知，但感兴趣的人很少；第二个系统中，收到通知的人很少，但这些人大多愿意参加。仅凭历史到场人数，无法区分两种解释。当新活动同时涉及消息转发、候补名单和日程重叠时，这种区别就会影响结果：转发决定哪些人知道机会，候补流转决定谁能获得名额，冲突则决定哪些参与意图最终转化为实际到场。这个例子同样适用于更广泛的资源分配与集体参与过程。校园是一个可执行实例，并不是研究问题的应用边界。

因此，从群体输出反推机制，是一个可能不可识别的逆问题。已有集合校准研究已经承认，多个参数配置可以同时满足统计上的相容要求 \citep{bai2022}。接下来的实际问题是：当逐条收集微观记录没有必要或代价过高时，应该补充哪一条过程观测？某条记录的行为熵很高，却可能对未来关心的预测几乎没有帮助。反过来，对于“是否查看过通知”的一个不大的分歧，如果能够区分未来拥挤轨迹不同的机制，就可能很有价值。目标导向实验设计提供了相应视角：证据价值取决于它要改善什么预测 \citep{chaloner1995,attia2018,zhong2026}。

本文用有限机制库、明确的环境结算和缓存化组合探针实现这一思路。共享参数控制注意、参与和候补退出；过程与轨迹分别接受相容性检查，形成候选集合。随后，轨迹加权规则定义一个校准混合分布，加权协方差按照“单位成本下预期减少多少探针预测分歧”对二元过程记录排序。同一缓存还支持后验更新与完整轨迹分布预测，查询步骤内部不需要重新训练模型或重跑模拟器。

实际结果比这一方法设计所暗示的更为有限。开发集选择强烈偏好集中的轨迹权重。最优预测器虽然在硬相容集合中保留了许多候选，但有效混合多样性几乎消失；轨迹分数改善的同时，个体过程的 Brier 分数却变差了。由于剩余可消除的机制分歧很小，取证消融的影响也不明显。这一现象揭示了一个直接影响方法有效性的张力：优化群体预测准确性，与为后续取证保留有用的不确定性，未必是同一个目标。

本文的贡献有三个方面。第一，针对严格区分行为意图与实际结果的多智能体模拟器，明确给出过程约束下的集合校准与预算化证据选择形式，并提供可复现实现。第二，将二元预测方差恒等式实现为批量协方差运算，在小批量中对假设标签积分，并说明计算与统计假设。我们不把该恒等式或目标导向设计原则本身当作新的理论贡献。第三，给出经过审计的合成研究，用同先验对照、逐活动不确定性和未知规则压力测试，区分校准收益与取证收益。结果表明，权重集中是一个实际失效模式，也约束了仅凭较好群体分数可以作出的推断。

本文应被理解为一项方法与诊断研究。它既没有证明在 SocioVerse 上达到当前最优，也没有取得 MIT 真实行为基准成绩。具有指定字段的接口不等于已经评测了相应数据集。一般数学形式可以迁移，但当前经验结论只针对这里的构造系统。
